\documentclass[11pt]{article}
\input{../preamble/preamble_latex.tex}
\input{../preamble/preamble_math.tex}


\title{STATS 305C: Week 5 Math Problem}
\date{Due Weds May 6, 2026 at 11:59pm}
\begin{document}
\maketitle
\noindent

Consider a hidden Markov model (HMM),
\begin{align*}
    z_1 &\sim \pi_0 \\
    z_t &\sim \pi_{z_{t-1}} && \text{for } t = 2, \ldots, T \\
    x_t &\sim p(x \mid \theta_{z_t}) && \text{for } t = 1, \ldots, T
\end{align*}
where $z_t \in \{1, \ldots, K\}$ denotes the discrete latent state at time $t$. The HMM
parameters consist of the initial state probabilities, $\pi_0$, the transition probabilities
$\{\pi_k\}_{k=1}^K$, and the emission parameters $\{\theta_k\}_{k=1}^K$.

\bigskip
\noindent
We can reparameterize the latent states in terms of their states at the times of change
points, $c_n \in \{1, \ldots, K\}$, and the corresponding durations between change points,
$d_n \in \mathbb{N}_+$ (these are positive integers). Let $c_1 = z_1$ and let $d_1$ denote
the number of time steps before the state changes. Then let $c_2 = z_{1+d_1}$ and let $d_2$
be the number of time steps before the state changes again. For example, this state sequence
would be reparameterized as follows,
\begin{align*}
    z_{1:T} = \underbrace{3, 3, 3, 3, 3,}_{c_1=3,\, d_1=5}
              \underbrace{1, 1, 1, 1,}_{c_2=1,\, d_2=4}
              \underbrace{2, 2, 2, 2, 2, 2,}_{c_3=2,\, d_3=6}
              \underbrace{1, 1, 1, 1, 1.}_{c_4=1,\, d_4=5}
\end{align*}

\begin{enumerate}[label=(\alph*)]
    \item Derive the probability mass function $p(d_n \mid c_n = k)$.

    \item Derive the probability mass function $p(c_{n+1} \mid c_n = k)$.

    \item Formulating the HMM in terms of states and durations suggests a natural extension
    to hidden \textit{semi-Markov} models (HSMMs), which allow for different distributions
    over the durations. For example, suppose we modeled the durations as following a negative
    binomial distribution,
    \begin{align*}
        p(d_n \mid c_n = k) = \mathrm{NB}(d_n - 1;\, r, q_k)
    \end{align*}
    where $r \in \mathbb{N}_+$ is a fixed positive integer and $q_k \in [0, 1]$ is a
    parameter. (The $d_n - 1$ on the r.h.s.\ comes from the fact that the negative binomial
    has support for all natural numbers including 0, whereas our durations must be at least
    1.) The negative binomial distribution is defined as the number of times a coin comes up
    tails before $r$ heads are observed, when $q_k$ is the probability of heads. How can an
    HSMM with negative binomial durations be reduced to an HMM on an extended state space?
\end{enumerate}

\end{document}